\chapter{Risultati}
\label{chap:risultati}

Questo capitolo presenta i risultati degli esperimenti definiti nei
\cref{chap:dataset,chap:modelli}. La domanda centrale è se il
bottleneck quantistico addestrabile Q migliori le prestazioni rispetto
al controllo classico B2, a parità delle condizioni di confronto
descritte nel protocollo. L'analisi considera sia la classificazione
binaria sia quella a tre classi e segue la variazione delle prestazioni
al crescere delle immagini disponibili per il training.

Si distingue tra prestazioni assolute, confronti appaiati Q--B2 e
analisi esplorative. Le prime descrivono il livello raggiunto dai
modelli; i secondi costituiscono il confronto statistico principale;
le analisi per classe, Q--QF, dimensionalità, calibrazione e costo
precisano la natura delle differenze osservate. Le possibili
spiegazioni dei risultati e le implicazioni più generali sono
riservate al capitolo di discussione.

\section{Completezza degli esperimenti e convenzioni di lettura}
\label{sec:ris-completezza}

La matrice finale comprende \(640\) run: \(160\) sul training completo
e \(160\) per ciascuno dei subset da \(100\), \(300\) e \(900\)
immagini. Per ogni regime sono presenti due task, otto configurazioni
e dieci repeat seed. Tutte le run risultano completate e dispongono
delle metriche di test. L'analisi non esclude seed sulla base del
livello di prestazione raggiunto.

Il protocollo è quello di repeated holdout descritto nel
\cref{chap:dataset}: per ciascun seed si usa la sola iterazione
selezionata, identificata negli artefatti da \texttt{iteration=1}.
Non si riportano risultati di una cross-validation completa su quattro
fold. Le dieci ripetizioni condividono il dataset di partenza e non
corrispondono a dieci campioni astronomici indipendenti. La lettura
dell'incertezza deve tenere conto della variabilità dovuta sia al
training sia al test, come evidenziato nella letteratura sul confronto
tra algoritmi di apprendimento \cite{nadeau2003inference}.

Le prestazioni provengono dai file \texttt{test\_metrics.json}; i
checkpoint erano stati selezionati sulla validation. Il controllo dei
risultati aggregati rispetto alle singole run ha confermato la
coincidenza delle Macro-F1, delle metriche probabilistiche controllate
e dei tempi. Il riepilogo contenuto nel pacchetto dei subset include
già anche le run Full: queste ultime sono contate una sola volta.

Nelle tabelle delle prestazioni assolute viene riportata
\(100\times\text{Macro-F1}\), mentre nei grafici resta la scala
\([0,1]\). Le differenze sono espresse in punti percentuali
(pp): una differenza di \(-0{,}0167\) nella Macro-F1 corrisponde
a circa \(-1{,}67\) pp, non a una diminuzione relativa dell'\(1{,}67\%\).
La notazione B2-4 o B2-8 indica la dimensione del bottleneck classico,
non la presenza di qubit in B2.

Gli effetti sono calcolati sui valori non arrotondati delle run.
La sottrazione dei valori già arrotondati nelle tabelle può quindi
differire di \(0{,}01\) pp dall'effetto riportato.

Le medie e gli intervalli di confidenza al \(95\%\) sono calcolati
a livello di repeat seed. Salvo indicazione diversa, gli intervalli
sono bootstrap percentile con \(10\,000\) ricampionamenti, secondo
il principio di ricampionamento di Efron \cite{efron1979bootstrap}. Le bande
delle figure rappresentano l'incertezza sulla media, non l'intera
dispersione delle singole run. I p-value riportati sono quelli dei
test appaiati archiviati, corretti con Holm \cite{holm1979multiple}
separatamente nella
famiglia primaria di quattro confronti Full e in quella secondaria
di dodici confronti sui subset. Gli intervalli rimangono marginali,
non simultanei. Restano valide le cautele inferenziali illustrate
nel \cref{chap:modelli}.

\section{Prestazioni assolute e numerosità del training}
\label{sec:ris-prestazioni}

\subsection{Classificazione binaria}

La \cref{tab:ris-binario} raccoglie tutte le configurazioni del task
binario. Sul Full le medie sono comprese tra \(96{,}00\) e
\(96{,}34\): le differenze assolute tra le architetture sono quindi
contenute. Anche i riferimenti B0 e B1 raggiungono circa \(96{,}33\)
e \(96{,}34\), senza che sia necessario introdurre un circuito
quantistico. Questi valori forniscono il contesto del problema,
ma non sostituiscono il confronto controllato Q--B2.

\newcommand{\risIC}[3]{\shortstack{#1\\{\scriptsize[#2; #3]}}}
\begin{table}[htbp]
\centering
\small
\setlength{\tabcolsep}{5pt}
\renewcommand{\arraystretch}{1.35}
\caption{Macro-F1 sul test binario, moltiplicata per 100: media e
IC95\% bootstrap sui dieci seed. Full corrisponde a 1\,876 immagini
di training. Le otto configurazioni sono tutte riportate, senza
selezione del miglior seed.}
\label{tab:ris-binario}
\begin{tabular}{lcccc}
\toprule
\rowcolor{ThesisTableHeader}
\tablehead{Modello} & \tablehead{100} & \tablehead{300} & \tablehead{900} & \tablehead{Full} \\
\midrule
B0 & \risIC{91,30}{90,55}{92,03} & \risIC{94,15}{93,41}{94,90} & \risIC{95,36}{94,86}{95,87} & \risIC{96,33}{95,90}{96,75} \\
B1 & \risIC{92,42}{91,39}{93,39} & \risIC{95,01}{94,51}{95,53} & \risIC{95,62}{95,29}{95,93} & \risIC{96,34}{96,00}{96,69} \\
\textbf{B2-4} & \risIC{92,50}{91,62}{93,45} & \risIC{94,55}{94,03}{95,04} & \risIC{95,94}{95,47}{96,32} & \risIC{96,22}{95,71}{96,67} \\
\textbf{B2-8} & \risIC{92,66}{91,82}{93,45} & \risIC{94,76}{94,49}{95,07} & \risIC{95,90}{95,40}{96,33} & \risIC{96,10}{95,46}{96,68} \\
\textbf{Q-4} & \risIC{90,90}{89,51}{92,05} & \risIC{94,16}{92,91}{95,27} & \risIC{95,96}{95,62}{96,30} & \risIC{96,24}{95,90}{96,57} \\
\textbf{Q-8} & \risIC{92,33}{91,02}{93,54} & \risIC{94,67}{93,80}{95,25} & \risIC{95,76}{95,13}{96,30} & \risIC{96,29}{95,84}{96,72} \\
QF-4 & \risIC{90,86}{89,80}{91,90} & \risIC{93,90}{92,70}{95,02} & \risIC{95,90}{95,54}{96,24} & \risIC{96,00}{95,49}{96,55} \\
QF-8 & \risIC{92,17}{90,77}{93,44} & \risIC{94,62}{93,75}{95,23} & \risIC{95,99}{95,47}{96,45} & \risIC{96,19}{95,67}{96,68} \\
\bottomrule
\end{tabular}
\end{table}

La riduzione dei dati rende più visibili alcune differenze. Con
\(100\) immagini, Q-4 raggiunge \(90{,}90\), contro \(92{,}50\)
di B2-4; Q-8 raggiunge invece \(92{,}33\), vicino al \(92{,}66\)
di B2-8. Nei regimi da \(900\) immagini e Full le medie di Q e
B2 si avvicinano per entrambe le dimensioni. Questa descrizione
non implica equivalenza statistica, che non è stata oggetto di un
test con margine prestabilito.

\subsection{Classificazione a tre classi}

Nel task a tre classi le prestazioni sono inferiori per tutte le
configurazioni e mostrano un aumento più ampio passando dai subset
ridotti al Full (\cref{tab:ris-treclassi}). B2-8 passa da
\(67{,}47\) a \(79{,}33\), mentre Q-8 passa da \(64{,}17\)
a \(77{,}66\). Il miglioramento descrittivo è rispettivamente di
\(11{,}86\) e \(13{,}49\) pp.

\begin{table}[htbp]
\centering
\small
\setlength{\tabcolsep}{5pt}
\renewcommand{\arraystretch}{1.35}
\caption{Macro-F1 sul test a tre classi, moltiplicata per 100:
media e IC95\% bootstrap sui dieci seed. Full corrisponde a
2\,814 immagini di training.}
\label{tab:ris-treclassi}
\begin{tabular}{lcccc}
\toprule
\rowcolor{ThesisTableHeader}
\tablehead{Modello} & \tablehead{100} & \tablehead{300} & \tablehead{900} & \tablehead{Full} \\
\midrule
B0 & \risIC{66,10}{64,84}{67,45} & \risIC{72,59}{71,42}{73,78} & \risIC{76,85}{75,90}{77,89} & \risIC{79,12}{77,99}{80,14} \\
B1 & \risIC{66,57}{64,88}{68,48} & \risIC{73,21}{72,16}{74,31} & \risIC{77,31}{76,47}{78,12} & \risIC{79,92}{78,98}{80,79} \\
\textbf{B2-4} & \risIC{67,44}{66,46}{68,38} & \risIC{73,36}{72,16}{74,64} & \risIC{76,59}{75,86}{77,41} & \risIC{79,00}{77,90}{80,06} \\
\textbf{B2-8} & \risIC{67,47}{66,20}{68,66} & \risIC{73,48}{72,51}{74,35} & \risIC{76,64}{75,88}{77,59} & \risIC{79,33}{78,34}{80,24} \\
\textbf{Q-4} & \risIC{64,55}{62,67}{66,45} & \risIC{71,02}{69,21}{72,73} & \risIC{75,74}{74,94}{76,48} & \risIC{78,27}{77,18}{79,20} \\
\textbf{Q-8} & \risIC{64,17}{62,90}{65,38} & \risIC{70,33}{69,26}{71,58} & \risIC{74,39}{72,45}{76,13} & \risIC{77,66}{76,86}{78,50} \\
QF-4 & \risIC{64,46}{62,42}{66,46} & \risIC{70,10}{68,27}{71,95} & \risIC{75,07}{74,10}{75,97} & \risIC{78,12}{77,08}{79,12} \\
QF-8 & \risIC{63,58}{61,91}{65,11} & \risIC{70,02}{68,87}{71,31} & \risIC{74,42}{72,45}{76,22} & \risIC{77,41}{76,42}{78,47} \\
\bottomrule
\end{tabular}
\end{table}

B2 mantiene una media superiore a Q in tutti e quattro i regimi,
sia nella configurazione a quattro dimensioni sia in quella a otto.
Anche B0 e B1 raggiungono sul Full valori superiori alle medie di
Q e QF. Non vengono però introdotti test a coppie tra tutte le
architetture: B0 e B1 rimangono riferimenti contestuali, con
architetture e condizioni di ottimizzazione non identiche a B2.

La \cref{fig:ris-scaling} mette a confronto i due task per le
configurazioni a otto dimensioni/qubit. Nel binario le traiettorie
sono vicine, soprattutto dai \(300\) campioni in poi. Nel task a
tre classi permane invece una separazione tra B2 e i modelli
quantistici. È importante precisarne l'andamento: aumentando i
dati, il divario medio Q--B2 si \emph{riduce}, passando da
\(-3{,}30\) pp con \(100\) immagini a \(-1{,}67\) pp sul Full;
non aumenta. L'incremento della numerosità migliora quindi entrambi
i modelli, senza annullare lo scarto nel protocollo osservato.

\begin{figure}[htbp]
\centering
% Titolo incorporato omesso per evitare ridondanza con la didascalia.
\includegraphics[width=\textwidth,trim=0 0 0 30bp,clip]{01_hero_learning_curves.pdf}
\caption{Prestazioni rispetto alla numerosità del training per B2-8,
Q-8 e QF-8. Punti e linee rappresentano le medie sui dieci seed;
le bande sono IC95\% bootstrap. Le posizioni 100, 300, 900 e Full
sono categoriche ed equidistanti: la pendenza dei segmenti non è
un tasso di miglioramento per immagine. I pannelli adottano scale
verticali diverse; Full contiene 1\,876 immagini nel binario e
2\,814 nel task a tre classi.}
\label{fig:ris-scaling}
\end{figure}

\section{Confronto appaiato Q--B2}
\label{sec:ris-qb2}

\subsection{Famiglia primaria: training completo}

Il confronto primario utilizza, per ogni seed, la differenza tra
la Macro-F1 di Q e quella di B2 con lo stesso task e la stessa
dimensione del bottleneck. Un valore negativo favorisce B2.
La \cref{tab:ris-primari} riporta tutti i quattro confronti Full,
con gli intervalli e i p-value conservati negli artefatti
dell'esperimento.

\begin{table}[htbp]
\centering
\small
\caption{Confronti primari Q--B2 sul Full. Dieci coppie per riga;
Holm è applicato alla famiglia dei quattro test. Gli IC95\% sono
bootstrap marginali. \(d_z\) è la differenza media divisa per la
deviazione standard delle differenze appaiate.}
\label{tab:ris-primari}
\begin{tabular}{lcrrrr}
\toprule
\rowcolor{ThesisTableHeader}
\tablehead{Task} & \tablehead{\(q\)} & \tablehead{\(\Delta\) (pp)} & \tablehead{IC95\% (pp)} & \tablehead{\(p_{\mathrm{Holm}}\)} & \tablehead{\(d_z\)} \\
\midrule
Binario & 4 & \(+0{,}02\) & \([-0{,}29;\,+0{,}33]\) & \(0{,}9406\) & \(+0{,}04\) \\
Binario & 8 & \(+0{,}19\) & \([-0{,}22;\,+0{,}61]\) & \(0{,}8797\) & \(+0{,}27\) \\
Tre classi & 4 & \(-0{,}73\) & \([-1{,}27;\,-0{,}33]\) & \(0{,}0132\) & \(-0{,}87\) \\
Tre classi & 8 & \(-1{,}67\) & \([-2{,}41;\,-1{,}02]\) & \(0{,}0135\) & \(-1{,}40\) \\
\bottomrule
\end{tabular}
\end{table}

Nel binario le differenze medie sono piccole e positive, ma gli
intervalli includono lo zero e nessun p-value corretto è inferiore
a \(0{,}05\). Non emerge quindi evidenza di un vantaggio di Q rispetto
a B2, né la mancanza di significatività consente di dichiararli
equivalenti.

Nel task a tre classi entrambi i confronti indicano invece una
Macro-F1 inferiore per Q, con \(p_{\mathrm{Holm}}<0{,}05\).
La perdita media è di \(0{,}73\) pp a quattro qubit e di
\(1{,}67\) pp a otto qubit. Gli intervalli sono interamente
negativi. La maggiore ampiezza osservata a otto qubit è una
descrizione degli effetti; non costituisce, da sola, un test della
differenza tra i due effetti.

\subsection{Famiglia secondaria: regimi low-resource}

La \cref{tab:ris-secondari} mostra i dodici confronti sui subset.
La correzione di Holm è calcolata su tutte le dodici righe, non
separatamente per task, dimensione o numerosità. Nel binario
nessun confronto risulta significativo dopo questa correzione.
Nel task a tre classi le differenze sono tutte negative; quelle
a otto qubit risultano significative in ciascun subset,
mentre quelle a quattro qubit non raggiungono la significatività.

\begin{table}[htbp]
\centering
\small
\caption{Tutti i confronti secondari Q--B2. Effetti e intervalli
sono espressi in punti percentuali. I p-value sono corretti con
Holm congiuntamente sui dodici confronti.}
\label{tab:ris-secondari}
\begin{tabular}{rlcrrr}
\toprule
\rowcolor{ThesisTableHeader}
\tablehead{Subset} & \tablehead{Task} & \tablehead{\(q\)} & \tablehead{\(\Delta\)} & \tablehead{IC95\%} & \tablehead{\(p_{\mathrm{Holm}}\)} \\
\midrule
100 & Binario & 4 & \(-1{,}61\) & \([-2{,}62;\,-0{,}56]\) & \(0{,}1632\) \\
100 & Binario & 8 & \(-0{,}33\) & \([-1{,}81;\,+0{,}82]\) & \(1{,}0000\) \\
100 & Tre classi & 4 & \(-2{,}88\) & \([-4{,}31;\,-1{,}29]\) & \(0{,}1071\) \\
100 & Tre classi & 8 & \(-3{,}30\) & \([-4{,}99;\,-1{,}80]\) & \(0{,}0396\) \\
\addlinespace
300 & Binario & 4 & \(-0{,}39\) & \([-1{,}42;\,+0{,}58]\) & \(1{,}0000\) \\
300 & Binario & 8 & \(-0{,}10\) & \([-0{,}81;\,+0{,}43]\) & \(1{,}0000\) \\
300 & Tre classi & 4 & \(-2{,}34\) & \([-3{,}76;\,-0{,}83]\) & \(0{,}1632\) \\
300 & Tre classi & 8 & \(-3{,}15\) & \([-4{,}46;\,-1{,}86]\) & \(0{,}0360\) \\
\addlinespace
900 & Binario & 4 & \(+0{,}02\) & \([-0{,}38;\,+0{,}40]\) & \(1{,}0000\) \\
900 & Binario & 8 & \(-0{,}14\) & \([-0{,}44;\,+0{,}11]\) & \(1{,}0000\) \\
900 & Tre classi & 4 & \(-0{,}85\) & \([-1{,}39;\,-0{,}25]\) & \(0{,}1632\) \\
900 & Tre classi & 8 & \(-2{,}24\) & \([-3{,}88;\,-0{,}94]\) & \(0{,}0396\) \\
\bottomrule
\end{tabular}
\end{table}

Con \(100\) immagini, il contrasto binario a quattro qubit è
\(-1{,}61\) pp e il suo IC95\% non include zero, ma il p-value
corretto è \(0{,}1632\). Non è una contraddizione: intervallo
marginale bootstrap e test corretto per la molteplicità non sono
lo stesso criterio decisionale. Lo stesso richiamo vale per i
contrasti a quattro qubit nel task a tre classi.

La \cref{fig:ris-delta} raccoglie i confronti nei quattro regimi.
Non si osserva un vantaggio statisticamente supportato di Q nei
regimi con meno dati. A otto qubit, nel task a tre classi, la
penalità media è anzi più ampia nei subset da \(100\) e \(300\)
immagini rispetto al Full. Si tratta di un andamento descrittivo
tra regimi: i test riportati valutano Q--B2 entro ciascun regime,
non la significatività della sua variazione tra numerosità.

\begin{figure}[htbp]
\centering
\includegraphics[width=\textwidth,trim=0 0 0 30bp,clip]{02_delta_q_minus_b2.pdf}
\caption{Differenze appaiate Q--B2 nei due task. Barre: differenza
media in pp; baffi: IC95\% bootstrap; asterisco:
\(p_{\mathrm{Holm}}<0{,}05\). I quattro confronti Full e i dodici
sui subset appartengono a famiglie di correzione distinte. Le
scale orizzontali dei pannelli sono diverse. Gli intervalli della
figura sono ricampionati dal notebook grafico e possono differire
lievemente da quelli archiviati nelle
\cref{tab:ris-primari,tab:ris-secondari}.}
\label{fig:ris-delta}
\end{figure}

\section{Passaggio dal binario alle tre classi}
\label{sec:ris-did}

La diminuzione della Macro-F1 assoluta interessa tutti i modelli.
Per isolare sul piano descrittivo la variazione del divario tra Q
e B2, si considera il contrasto
\[
 D_s^{(\mathrm{pp})}=100D_s
 =100\left[(F^{\mathrm Q}_{1,s}-F^{\mathrm{B2}}_{1,s})_{\mathrm{tre}}
 -(F^{\mathrm Q}_{1,s}-F^{\mathrm{B2}}_{1,s})_{\mathrm{binario}}\right].
\]
Qui \(D_s\) è definito nella \cref{eq:did} sulla scala originale
della Macro-F1. Le differenze sono costruite entro lo stesso repeat seed; il
bootstrap ricampiona questi dieci contrasti, non quattro medie
indipendenti. Un valore negativo indica che Q perde terreno
rispetto a B2 nel task a tre classi.

\begin{table}[htbp]
\centering
\small
\caption{Differenza-in-differenze tra task, in pp: media e
IC95\% bootstrap dei contrasti appaiati. Analisi esplorativa,
senza un'ulteriore famiglia di test o correzioni di Holm.}
\label{tab:ris-did}
\begin{tabular}{lrr}
\toprule
\rowcolor{ThesisTableHeader}
\tablehead{Regime} & \tablehead{4 qubit: media [IC95\%]} & \tablehead{8 qubit: media [IC95\%]} \\
\midrule
100 & \(-1{,}27\;[-3{,}13;\,+0{,}86]\) & \(-2{,}97\;[-5{,}17;\,-0{,}78]\) \\
300 & \(-1{,}95\;[-3{,}37;\,-0{,}46]\) & \(-3{,}06\;[-4{,}61;\,-1{,}49]\) \\
900 & \(-0{,}87\;[-1{,}59;\,-0{,}02]\) & \(-2{,}10\;[-3{,}54;\,-0{,}91]\) \\
Full & \(-0{,}75\;[-1{,}42;\,-0{,}24]\) & \(-1{,}86\;[-2{,}56;\,-1{,}21]\) \\
\bottomrule
\end{tabular}
\end{table}

Tutte le stime medie sono negative (\cref{tab:ris-did}). A otto
qubit il contrasto è compreso tra circa \(-3{,}06\) e
\(-1{,}86\) pp e gli intervalli marginali restano sotto lo zero.
A quattro qubit gli effetti sono più contenuti; con \(100\)
immagini l'intervallo attraversa lo zero. Il limite superiore
molto vicino a zero per \(900\) immagini non va interpretato
come una separazione netta tra presenza e assenza dell'effetto.

Questa analisi rende visibile una differenza relativa che sarebbe
poco leggibile rappresentando soltanto il calo comune di circa
venti punti della Macro-F1. Non identifica però un effetto
causale dell'aggiunta della classe merger: cambiano il problema
di classificazione, la composizione dei dati e, sul Full, anche
la numerosità totale del training. La Macro-F1 dei due task è
inoltre una media su un diverso numero di classi.

\section{Analisi della classe merger}
\label{sec:ris-merger}

\subsection{F1, precision e recall}

La classe merger è quella con la F1 media più bassa sul Full
per tutte le configurazioni B2, Q e QF riportate nella
\cref{tab:ris-merger}. Il confronto con le altre due classi
permette di localizzare il divario, senza attribuirgli
automaticamente una causa fisica o un problema nelle etichette.

\begin{table}[htbp]
\centering
\small
\caption{Prestazioni medie per classe sul Full a tre classi,
moltiplicate per 100. Per merger sono riportate anche precision
e recall. Ogni valore è la media della metrica calcolata
separatamente sui dieci test; F1 non è ricostruita dalla
precision e dal recall medi.}
\label{tab:ris-merger}
\begin{tabular}{lrrrrr}
\toprule
& \multicolumn{3}{c}{\tablehead{F1 per classe}} & \multicolumn{2}{c}{\tablehead{Merger}} \\
\cmidrule(lr){2-4}\cmidrule(lr){5-6}
\rowcolor{ThesisTableHeader}
\tablehead{Modello} & \tablehead{Ellittiche} & \tablehead{Spirali} & \tablehead{Merger} & \tablehead{Precision} & \tablehead{Recall} \\
\midrule
\textbf{B2-4} & 84,99 & 81,76 & 70,27 & 72,68 & 68,23 \\
Q-4  & 84,69 & 80,74 & 69,38 & 71,04 & 68,53 \\
QF-4 & 84,70 & 80,49 & 69,17 & 71,59 & 67,59 \\
\addlinespace
\textbf{B2-8} & 85,16 & 82,26 & 70,58 & 73,79 & 67,87 \\
Q-8  & 84,61 & 80,39 & 67,99 & 71,63 & 64,97 \\
QF-8 & 84,83 & 79,66 & 67,73 & 71,53 & 64,90 \\
\bottomrule
\end{tabular}
\end{table}

A otto qubit, Q perde rispetto a B2 circa \(2{,}58\) pp nella
F1 merger, \(1{,}86\) pp nelle spirali e \(0{,}56\) pp nelle
ellittiche. Poiché la Macro-F1 è la media delle tre F1, il
contributo aritmetico della classe merger al divario complessivo
è circa \(2{,}58/3=0{,}86\) pp, poco più della metà dei
\(1{,}67\) pp osservati. Questo è un modo di scomporre la metrica,
non una spiegazione causale degli errori.

A quattro qubit il quadro è meno concentrato sulla merger:
la diminuzione della sua F1 è di \(0{,}88\) pp, mentre per le
spirali è di \(1{,}02\) pp. Inoltre il recall merger di Q-4 è
lievemente superiore a quello di B2-4 (\(68{,}53\) contro
\(68{,}23\)), a fronte di una precision inferiore
(\(71{,}04\) contro \(72{,}68\)). Non sarebbe quindi corretto
descrivere entrambe le dimensioni come un'identica perdita di
capacità di recuperare le merger.

\subsection{Andamento nei subset ed errori di classificazione}

La \cref{fig:ris-merger} mostra che il divario sulla F1 merger
è particolarmente ampio con \(100\) immagini e otto qubit:
Q--B2 è pari a \(-9{,}82\) pp, con IC95\% esplorativo
\([-13{,}32;\,-6{,}72]\). La differenza si riduce a
\(-4{,}82\) pp con \(300\) immagini, \(-3{,}50\) pp con
\(900\) e \(-2{,}58\) pp sul Full. A quattro qubit il divario
medio passa da \(-4{,}47\) pp nel subset minimo a
\(-0{,}88\) pp sul Full, ma non decresce monotonamente:
con \(900\) immagini è di soli \(-0{,}32\) pp.

\begin{figure}[htbp]
\centering
\includegraphics[width=\textwidth,trim=0 0 0 30bp,clip]{04_merger_detailed.pdf}
\caption{F1 della classe merger rispetto al regime di training,
con pannelli separati per quattro e otto dimensioni/qubit.
Le linee rappresentano le medie e le bande gli IC95\% bootstrap.
Le due scale verticali coincidono. La figura descrive la F1 di
una singola classe, non la Macro-F1 dell'intero task.}
\label{fig:ris-merger}
\end{figure}

Le matrici di confusione normalizzate forniscono un riscontro
più specifico sugli errori del Full. Per ogni modello sono state
mediate le matrici normalizzate per riga dei dieci seed, senza
trattare come oggetti distinti le ricomparse di una stessa
galassia in holdout diversi.
In B2-8, in media, il \(17{,}44\%\) delle merger è classificato
come spirale; in Q-8 la quota sale al \(20{,}68\%\).
La quota di merger assegnate alle ellittiche è invece simile:
\(14{,}69\%\) in B2-8 e \(14{,}35\%\) in Q-8.
La diminuzione del recall merger a otto qubit si associa dunque,
in questi tassi medi, soprattutto a più assegnazioni alla classe
spirale.

A quattro qubit, invece, le merger assegnate alle spirali sono
il \(18{,}78\%\) per B2-4 e il \(17{,}91\%\) per Q-4.
Aumenta la confusione nella direzione opposta: le spirali
assegnate alla classe merger passano dal \(13{,}69\%\) al
\(16{,}12\%\). Questo andamento è coerente con la minore
precision merger di Q-4. Le percentuali descrivono tassi medi
per classe vera, non percentuali sull'intero dataset.

\section[Addestramento quantistico e dimensionalità]{Contributo dell'addestramento quantistico e dimensionalità}
\label{sec:ris-ablazioni}

\subsection{Circuito addestrabile rispetto al circuito congelato}

Il confronto Q--QF valuta quanto cambi la prestazione quando i
pesi del circuito vengono ottimizzati, mantenendo addestrabili
le componenti classiche in entrambi i modelli. QF non è quindi
un classificatore interamente congelato.

Sul Full il vantaggio medio di Q è contenuto: nel binario è di
\(0{,}23\) pp a quattro qubit e \(0{,}10\) pp a otto;
nel task a tre classi è di \(0{,}15\) e \(0{,}26\) pp.
Tutti i quattro intervalli bootstrap Full includono lo zero
(\cref{fig:ris-qqf}). Questi dati non documentano un beneficio
netto e uniforme dell'ottimizzazione dei pesi circuitali.

\begin{figure}[htbp]
\centering
\includegraphics[width=\textwidth,trim=0 0 0 30bp,clip]{05_q_minus_qf.pdf}
\caption{Confronto esplorativo Q--QF: differenza media appaiata
di Macro-F1 in pp e IC95\% bootstrap. Valori positivi favoriscono
Q. I colori indicano il segno dell'effetto, non la significatività.
Le scale orizzontali dei due task sono diverse. Non sono
introdotti nuovi p-value per questi confronti.}
\label{fig:ris-qqf}
\end{figure}

Alcuni subset mostrano effetti più favorevoli a Q. Nel task a
tre classi a quattro qubit, Q--QF è \(+0{,}92\) pp con
\(300\) immagini, con IC95\% circa \([+0{,}24;\,+1{,}70]\),
e \(+0{,}67\) pp con \(900\), con IC95\% circa
\([+0{,}21;\,+1{,}30]\). Nel binario a otto qubit e
\(100\) immagini il vantaggio è più piccolo, \(+0{,}15\) pp,
con IC95\% circa \([+0{,}04;\,+0{,}28]\). Al contrario,
Q-8 nel binario con \(900\) immagini è mediamente inferiore
a QF-8 di \(0{,}24\) pp, con intervallo che include lo zero.
Sono evidenze esplorative e non una nuova serie di risultati
confermativi selezionati in base agli intervalli più favorevoli.

\subsection{Da quattro a otto dimensioni/qubit}

La sensibilità alla dimensione è riassunta nella
\cref{fig:ris-qubit}. Nel binario con \(100\) immagini,
Q-8 supera Q-4 di \(1{,}43\) pp; sul Full la differenza si
riduce a \(0{,}05\) pp. Nel task a tre classi, invece, Q-8
ha una media inferiore a Q-4 in tutti i regimi: le differenze
\(8-4\) sono circa \(-0{,}38\), \(-0{,}69\), \(-1{,}34\)
e \(-0{,}61\) pp per \(100\), \(300\), \(900\) e Full.
Per B2, nello stesso task, le differenze sono piccole e positive.

\begin{figure}[htbp]
\centering
\includegraphics[width=\textwidth,trim=0 0 0 30bp,clip]{06_qubit_sensitivity.pdf}
\caption{Variazione media della Macro-F1 passando da quattro
a otto dimensioni/qubit, in pp. Valori positivi favoriscono
la dimensione otto. La heatmap è una sintesi descrittiva senza
intervalli o marcature di significatività. Per B2 varia la
dimensione del bottleneck classico; per Q e QF cambiano anche
il circuito e le dimensioni dei livelli classici collegati.}
\label{fig:ris-qubit}
\end{figure}

L'aumento della dimensione non coincide dunque con un incremento
sistematico delle prestazioni. Questa osservazione riguarda le
configurazioni effettivamente confrontate: come precisato nel
\cref{chap:modelli}, il passaggio da quattro a otto qubit modifica
anche la proiezione classica e il classificatore finale. Non
isola l'effetto del solo numero di qubit.

\section{Qualità delle probabilità e calibrazione}
\label{sec:ris-calibrazione}

La Macro-F1 valuta le decisioni di classe, ma non esaurisce
l'analisi delle probabilità prodotte dal modello. La
\cref{tab:ris-calibrazione} riporta ECE, Brier score e log-loss
medi sul Full per B2 e Q. Valori più bassi sono preferibili,
ma le tre quantità misurano aspetti diversi; in particolare
Brier score e log-loss non sono misure di sola calibrazione
\cite{gneiting2007scoring}. L'ECE confronta invece accuratezza e
confidenza aggregate in intervalli \cite{guo2017calibration}.

\begin{table}[htbp]
\centering
\small
\caption{Metriche probabilistiche medie sul test Full, senza
ricalibrazione post-hoc. Dieci seed per configurazione. I valori
sono nella scala originale delle metriche, non moltiplicati
per 100. Sintesi descrittiva, senza test aggiuntivi.}
\label{tab:ris-calibrazione}
\begin{tabular}{llrrr}
\toprule
\rowcolor{ThesisTableHeader}
\tablehead{Task} & \tablehead{Modello} & \tablehead{ECE} & \tablehead{Brier score} & \tablehead{Log-loss} \\
\midrule
Binario & B2-4 & 0,0189 & 0,0603 & 0,1205 \\
        & Q-4  & 0,0158 & 0,0610 & 0,1172 \\
        & B2-8 & 0,0205 & 0,0631 & 0,1258 \\
        & Q-8  & 0,0158 & 0,0598 & 0,1160 \\
\addlinespace
Tre classi & B2-4 & 0,0495 & 0,3024 & 0,5137 \\
           & Q-4  & 0,0508 & 0,3132 & 0,5361 \\
           & B2-8 & 0,0570 & 0,2970 & 0,5078 \\
           & Q-8  & 0,0432 & 0,3185 & 0,5417 \\
\bottomrule
\end{tabular}
\end{table}

Nel binario, Q mostra ECE e log-loss medi inferiori a B2 per
entrambe le dimensioni; il Brier score è lievemente superiore
per Q-4 e inferiore per Q-8. Nel task a tre classi, Brier score
e log-loss sono maggiori per Q in entrambe le configurazioni.
L'ECE non segue uniformemente questo ordinamento: Q-8 ha un
ECE medio di \(0{,}0432\), inferiore al \(0{,}0570\) di B2-8,
pur avendo una Macro-F1 inferiore e peggiori Brier score e
log-loss.

Non vi è contraddizione tra questi risultati: una minore
discrepanza media tra confidenza e accuratezza nei bin dell'ECE
non garantisce una migliore discriminazione delle classi o una
migliore qualità probabilistica secondo ogni criterio. Non si
utilizza quindi il solo ECE per ribaltare il confronto sulla
Macro-F1, né un reliability diagram di un singolo seed per
caratterizzare un'intera architettura. I valori assoluti delle
metriche probabilistiche tra binario e tre classi non vanno
letti come se composizione e numero di classi fossero invariati.

\section{Prestazioni e costo computazionale}
\label{sec:ris-costi}

\subsection{Tempi registrati e ambito del confronto}

I tempi sono confrontati tra modelli dello stesso regime, secondo
le condizioni di esecuzione descritte nella \cref{sec:ambiente-hardware}.
Le traiettorie tra regimi hanno funzione descrittiva e non misurano
lo scaling del costo rispetto alla sola numerosità del training.

La \cref{tab:ris-tempi-full} riporta i tempi mediani del Full.
Il tempo di training è la somma dei tempi di epoca registrati:
non include estrazione iniziale delle feature, salvataggio
successivo dei checkpoint, valutazione finale o trasferimenti
degli archivi. Il tempo di predizione qui riportato è invece
lo specifico campo \texttt{test\_inference\_seconds}, riferito
all'intero test dopo il ripristino del checkpoint selezionato.
Non è il campo \texttt{inference\_seconds} delle tabelle
intermedie, che si riferisce alla validation.

\begin{table}[htbp]
\centering
\small
\caption{Tempi mediani sul Full in secondi, sui dieci seed.
Training: somma dei tempi di epoca; test: tempo di predizione
su 938 esempi nel binario e 1\,407 nel task a tre classi.
Le misure escludono la costruzione della cache di feature e
non rappresentano latenze end-to-end delle immagini.}
\label{tab:ris-tempi-full}
\begin{tabular}{lrrrr}
\toprule
& \multicolumn{2}{c}{\tablehead{Binario}} & \multicolumn{2}{c}{\tablehead{Tre classi}} \\
\cmidrule(lr){2-3}\cmidrule(lr){4-5}
\rowcolor{ThesisTableHeader}
\tablehead{Modello} & \tablehead{Training} & \tablehead{Test} & \tablehead{Training} & \tablehead{Test} \\
\midrule
\textbf{B2-4} & 23,88 & 0,192 & 41,90 & 0,287 \\
Q-4  & 250,27 & 2,253 & 441,80 & 3,456 \\
QF-4 & 209,17 & 2,188 & 380,25 & 3,273 \\
\textbf{B2-8} & 23,31 & 0,195 & 42,41 & 0,279 \\
Q-8  & 442,67 & 4,211 & 854,89 & 6,557 \\
QF-8 & 361,37 & 4,181 & 668,88 & 6,692 \\
\bottomrule
\end{tabular}
\end{table}

I tempi di Q sono maggiori di quelli di B2 sia durante il
training sia durante la predizione sul test. QF riduce il
tempo mediano di training rispetto a Q nelle configurazioni
Full, ma resta molto più costoso del controllo classico.
Il congelamento dei pesi circuitali non elimina infatti
l'esecuzione del circuito né la propagazione del gradiente
necessaria ad addestrare la proiezione che lo precede.

\subsection{Rapporti appaiati e compromesso prestazioni--tempo}

La \cref{tab:ris-rapporti} riporta la mediana dei dieci rapporti
\(t_{\mathrm Q,s}/t_{\mathrm{B2},s}\), non il rapporto tra i
tempi mediani della tabella precedente. I rapporti mediani
di training vanno da circa \(10\) a \(15\) a quattro qubit
e da circa \(19\) a \(26\) a otto qubit.

\begin{table}[htbp]
\centering
\small
\caption{Mediana dei rapporti appaiati di tempo di training
Q/B2. Un valore 10 indica un tempo dieci volte quello di B2
secondo questa statistica. Le colonne dei subset provengono
da CPU, la colonna Full da GPU; non sono fattori di scaling
misurati a hardware costante.}
\label{tab:ris-rapporti}
\begin{tabular}{lcrrrr}
\toprule
\rowcolor{ThesisTableHeader}
\tablehead{Task} & \tablehead{\(q\)} & \tablehead{100} & \tablehead{300} & \tablehead{900} & \tablehead{Full} \\
\midrule
Binario & 4 & 14,63 & 11,03 & 10,15 & 10,19 \\
Binario & 8 & 25,66 & 19,10 & 18,97 & 19,26 \\
Tre classi & 4 & 12,18 & 11,14 & 11,11 & 10,58 \\
Tre classi & 8 & 21,44 & 21,56 & 23,95 & 18,58 \\
\bottomrule
\end{tabular}
\end{table}

La \cref{fig:ris-costo} combina Macro-F1 media e tempo di
training mediano. Nel task a tre classi, entro ogni regime,
B2 raggiunge una prestazione media maggiore con un tempo
mediano inferiore rispetto alle corrispondenti configurazioni
Q e QF. Nel binario alcuni punti quantistici hanno una
Macro-F1 media lievemente superiore, ma a fronte di un
costo temporale nettamente maggiore. Non è dunque corretto
descrivere B2 come strettamente dominante in ogni singolo
confronto binario soltanto sulla base del grafico.

\begin{figure}[htbp]
\centering
\includegraphics[width=\textwidth,trim=0 0 0 30bp,clip]{07_pareto_performance_cost.pdf}
\caption{Prestazioni e tempo di training osservato: ciascun
punto combina la Macro-F1 media e il tempo mediano delle
dieci run di una configurazione. Le linee collegano, in
ordine, 100, 300, 900 e Full; le etichette identificano
l'estremo Full. L'asse dei tempi è logaritmico. I primi
tre regimi sono eseguiti su CPU, Full su GPU: i segmenti
non definiscono una frontiera a hardware costante né una
legge di scaling. La figura non rappresenta l'incertezza
congiunta di prestazione e tempo.}
\label{fig:ris-costo}
\end{figure}

L'early stopping contribuisce ai tempi complessivi, che
dipendono anche dal numero di epoche effettivamente eseguite.
Le differenze non possono essere attribuite interamente
al costo di una singola valutazione circuitale. Infine,
tutti questi risultati riguardano il simulatore classico
\texttt{default.qubit}: non misurano tempi di una QPU,
costi di accesso all'hardware quantistico o consumi energetici.

\section{Sintesi delle evidenze}
\label{sec:ris-sintesi}

Nel protocollo considerato, le prestazioni medie di tutte le
configurazioni crescono all'aumentare dei dati di training.
Sul Full binario i modelli raggiungono Macro-F1 vicine a
\(0{,}96\), senza evidenza di un vantaggio di Q rispetto
a B2 nei confronti primari. Nel task a tre classi, invece,
Q presenta una Macro-F1 inferiore a B2 per entrambe le
dimensioni, con significatività secondo la procedura
appaiata e la correzione adottate.

Nei regimi low-resource non emerge il vantaggio quantistico
ipotizzato: a otto qubit, nel task a tre classi, le perdite
medie sono maggiori di quelle osservate sul Full e risultano
significative nella famiglia secondaria. L'analisi merger
localizza una parte importante del divario, soprattutto
a otto qubit, senza spiegare da sola tutti gli errori.
Il confronto con QF mostra benefici dell'addestramento
circuitale piccoli sul Full e dipendenti dal regime nei
subset. Aumentare da quattro a otto qubit non offre un
miglioramento sistematico della Macro-F1.

Le metriche probabilistiche aggiungono un quadro non
riducibile al solo ordinamento della Macro-F1; i costi
registrati mostrano invece un aumento consistente del
tempo richiesto dai modelli quantistici simulati rispetto
a B2. Queste evidenze circoscrivono il risultato sperimentale
alle architetture, ai dati e alle condizioni qui analizzati:
non costituiscono un'affermazione generale sull'impossibilità
di un vantaggio del machine learning quantistico.
